\documentclass[12pt,a4paper]{article}

% ===== Core Packages =====
\usepackage[utf8]{inputenc}
\usepackage[T5]{fontenc}
\usepackage[vietnamese,english]{babel}
\usepackage{amsmath, amssymb, amsthm}
\usepackage{mathtools}
\usepackage{algorithm}
\usepackage{algpseudocode}
\usepackage{graphicx}
\usepackage{xcolor}
\usepackage{booktabs}
\usepackage{array}
\usepackage{enumitem}
\usepackage{hyperref}
\usepackage{geometry}
\usepackage{tcolorbox}
\usepackage{tikz}
\usepackage{pgfplots}
\usepackage{caption}
\usepackage{subcaption}
\usepackage{bm}
\usepackage{bbm}
\usepackage{lmodern}
\usepackage{fontawesome5}
\usepackage{microtype}

\geometry{margin=2.5cm}
\pgfplotsset{compat=1.18}
\usetikzlibrary{arrows.meta, positioning, shapes.geometric}

% ===== Colors =====
\definecolor{primary}{RGB}{0,90,170}
\definecolor{accent}{RGB}{200,50,60}
\definecolor{evgreen}{RGB}{0,130,80}
\definecolor{rideorange}{RGB}{200,90,0}
\definecolor{boxbg}{RGB}{235,245,255}
\definecolor{formulabg}{RGB}{255,250,235}
\definecolor{evbg}{RGB}{232,248,240}
\definecolor{ridebg}{RGB}{255,245,230}

% ===== Custom tcolorboxes =====
\tcbuselibrary{skins, breakable, theorems}

\newtcolorbox{definitionbox}[1]{
  colback=boxbg, colframe=primary,
  title=\textbf{#1}, fonttitle=\color{white}\bfseries,
  colbacktitle=primary, breakable, arc=4pt
}
\newtcolorbox{formulabox}[1]{
  colback=formulabg, colframe=accent,
  title=\textbf{#1}, fonttitle=\color{white}\bfseries,
  colbacktitle=accent, breakable, arc=4pt
}
\newtcolorbox{evbox}[1]{
  colback=evbg, colframe=evgreen,
  title=\textbf{#1}, fonttitle=\color{white}\bfseries,
  colbacktitle=evgreen, breakable, arc=4pt
}
\newtcolorbox{ridebox}[1]{
  colback=ridebg, colframe=rideorange,
  title=\textbf{#1}, fonttitle=\color{white}\bfseries,
  colbacktitle=rideorange, breakable, arc=4pt
}

% ===== Theorem environments =====
\newtheorem{theorem}{Định lý}[section]
\newtheorem{proposition}{Mệnh đề}[section]
\newtheorem{definition}{Định nghĩa}[section]
\newtheorem{remark}{Nhận xét}[section]

% ===== Math operators =====
\DeclareMathOperator{\argmin}{arg\,min}
\DeclareMathOperator{\argmax}{arg\,max}
\DeclareMathOperator*{\E}{\mathbb{E}}
\DeclareMathOperator{\Prob}{\mathbb{P}}
\DeclareMathOperator{\CVaR}{CVaR}
\DeclareMathOperator{\Gini}{Gini}

\hypersetup{colorlinks=true, linkcolor=primary, urlcolor=primary, citecolor=primary}

\title{
  \textbf{\Large Tối ưu hóa Dịch vụ Xe \& Giao hàng}\\[0.4em]
  \large Công thức Toán học, Mô hình DS/AI và Thuật toán\\[0.3em]
  \normalsize Tài liệu Kỹ thuật Chuyên sâu --- Green SM / Vingroup LogTech
}
\author{Phiên bản 2.0 --- Tháng 5, 2026}
\date{}

%% ================================================================
\begin{document}
\maketitle
\tableofcontents
\newpage

%% ================================================================
\section{Tổng quan về Pipeline Tối ưu hóa}

\subsection{Phân tầng quyết định}

Trong hệ sinh thái LogTech của Green SM và Vingroup, các bài toán tối ưu hóa
vận hành tạo thành một \emph{chuỗi phân tầng} (hierarchical optimization pipeline):

\begin{itemize}[leftmargin=2em, itemsep=4pt]
  \item \textbf{Tầng Chiến lược (Strategic):} Quy hoạch mạng lưới trạm sạc,
        kho trung chuyển, phân bổ đội xe điện dị chủng --- cập nhật theo ngày/tuần.
  \item \textbf{Tầng Chiến thuật (Tactical):} Dự báo cầu, điều phối tài xế,
        định giá surge, quản lý SoC đội xe --- cập nhật theo giờ.
  \item \textbf{Tầng Vận hành (Operational):} Ghép đơn, định tuyến, dispatch,
        lịch sạc real-time --- độ trễ mili-giây đến giây.
\end{itemize}

\subsection{Ký hiệu chung}

\begin{table}[h!]
\centering
\caption{Ký hiệu toán học sử dụng xuyên suốt tài liệu}
\renewcommand{\arraystretch}{1.25}
\begin{tabular}{>{\ttfamily}ll}
\toprule
\textbf{Ký hiệu} & \textbf{Ý nghĩa} \\
\midrule
$\mathcal{O} = \{o_1,\ldots,o_n\}$ & Tập $n$ đơn hàng / chuyến khách \\
$\mathcal{D} = \{d_1,\ldots,d_m\}$ & Tập $m$ tài xế / phương tiện \\
$\mathcal{G} = (\mathcal{V},\mathcal{E})$ & Đồ thị đường đi \\
$\mathcal{F}$ & Tập trạm sạc điện \\
$\tau(u,v)$ & Thời gian di chuyển từ $u$ đến $v$ \\
$c_{ij}$ & Chi phí từ node $i$ đến node $j$ \\
$[a_i,b_i]$ & Cửa sổ thời gian của node $i$ \\
$q_i$ & Tải / số hành khách của đơn $i$ \\
$\mathbf{Q}_k$ & Vector năng lực xe $k$ (xem Mục~\ref{sec:hetfleet}) \\
$e_k(t)$ & State-of-Charge (SoC) xe $k$ tại thời điểm $t$ \\
$\pi$ & Policy trong Reinforcement Learning \\
\bottomrule
\end{tabular}
\end{table}

%% ================================================================
\section{Bài toán Ghép đơn --- Order Batching \& Pooling}

\subsection{Phát biểu bài toán}

\begin{definitionbox}{Bài toán Ghép đơn (Order Batching Problem --- OBP)}
Tìm cách phân chia $\mathcal{O}$ thành các \emph{batch}
$\mathcal{B} = \{B_1,\ldots,B_K\}$ sao cho
$\mathcal{O} = B_1 \cup \cdots \cup B_K$, $B_i \cap B_j = \emptyset$,
mỗi batch gán cho một tài xế, tối thiểu hóa tổng chi phí hệ thống.
\end{definitionbox}

\subsection{Mô hình ILP (đã linearize)}

\subsubsection{Biến quyết định}

\begin{alignat}{2}
  y_{ik} &\in \{0,1\} \quad&&\text{: Đơn $i$ được gán cho tài xế $k$}\\
  z_k    &\in \{0,1\} \quad&&\text{: Tài xế $k$ được kích hoạt}\\
  w_{ijk}&\in \{0,1\} \quad&&\text{: Đơn $i$ và $j$ cùng giao cho tài xế $k$ (linearization)}
\end{alignat}

\subsubsection{Hàm mục tiêu}

\begin{formulabox}{Hàm mục tiêu OBP --- dạng ILP tuyến tính}
\begin{equation}
  \min \;
  \underbrace{\sum_{k}\sum_{i<j} c_{ij}\,w_{ijk}}_{\text{Chi phí di chuyển (tuyến tính)}}
  + \lambda\underbrace{\sum_i \max\!\bigl(0,\,t_i^{\text{deliver}} - b_i\bigr)}_{\text{Phạt trễ}}
  + \mu\underbrace{\sum_k z_k}_{\text{Kích hoạt}}
\end{equation}
\end{formulabox}

\begin{remark}[McCormick Linearization]
Hàm gốc $\sum c_{ij}y_{ik}y_{jk}$ là \textbf{QIP} (bậc 2). Ta đặt
$w_{ijk} = y_{ik}y_{jk}$ và tuyến hóa bằng \emph{McCormick envelopes}:
\begin{equation}
  w_{ijk} \leq y_{ik},\quad w_{ijk} \leq y_{jk},\quad
  w_{ijk} \geq y_{ik} + y_{jk} - 1,\quad w_{ijk} \geq 0.
\end{equation}
\end{remark}

\subsubsection{Ràng buộc}

\begin{alignat}{3}
  &\sum_k y_{ik} = 1 &&\forall i &&\quad\text{(C1: Đúng 1 tài xế)}\\
  &\sum_i q_i\,y_{ik} \leq Q_k\,z_k &&\forall k &&\quad\text{(C2: Sức chứa)}\\
  &t_i^{\text{deliver}} \geq a_i &&\forall i &&\quad\text{(C3: Time window)}\\
  &|B_k| \leq B_{\max} &&\forall k &&\quad\text{(C4: Giới hạn batch)}\\
  &w_{ijk} \leq y_{ik},\; w_{ijk} \leq y_{jk} &&\forall i,j,k &&\quad\text{(C5: McCormick upper)}\\
  &w_{ijk} \geq y_{ik}+y_{jk}-1 &&\forall i,j,k &&\quad\text{(C6: McCormick lower)}
\end{alignat}

\subsection{Phân cụm không-thời gian (Spatiotemporal Clustering)}

\begin{formulabox}{Khoảng cách không-thời gian (Haversine)}
\begin{equation}
  d(o_i,o_j)
  = \alpha\,d_{\text{geo}}^{pu}(i,j)
  + \beta\,d_{\text{geo}}^{do}(i,j)
  + \gamma\,|t_i^{req}-t_j^{req}|
\end{equation}
\begin{equation}
  d_{\text{geo}}(i,j)
  = 2R\arcsin\sqrt{\sin^2\!\tfrac{\Delta\phi}{2}
    + \cos\phi_i\cos\phi_j\sin^2\!\tfrac{\Delta\lambda}{2}},\quad R = 6371\text{ km}
\end{equation}
\end{formulabox}

\begin{algorithm}[H]
\caption{ST-DBSCAN cho Ghép đơn cùng tuyến}
\begin{algorithmic}[1]
\Require $\mathcal{O},\;\varepsilon_s,\;\varepsilon_t,\;\mathrm{MinPts}$
\Ensure  $\mathcal{B}=\{B_1,\ldots,B_K\}$
\State Khởi tạo: mọi đơn là \textsc{unvisited}
\For{mỗi $o_i\in\mathcal{O}$}
  \If{$o_i$ là \textsc{unvisited}}
    \State $N\leftarrow\mathrm{RangeQuery}(o_i,\varepsilon_s,\varepsilon_t)$
    \If{$|N|<\mathrm{MinPts}$} Đánh dấu noise
    \Else\; Tạo $B_k$; mở rộng cụm từ $N$ \EndIf
  \EndIf
\EndFor
\State \Return $\mathcal{B}$
\end{algorithmic}
\end{algorithm}

\begin{algorithm}[H]
\caption{Greedy Insertion --- thứ tự Pickup-Delivery}
\begin{algorithmic}[1]
\Require Batch $B_k$, vị trí tài xế $\ell_k$
\Ensure  Lộ trình $\pi$
\State $\pi\leftarrow[\;]$; unvisited $\leftarrow\bigcup_i\{pu_i,do_i\}$
\While{unvisited $\neq\emptyset$}
  \State $v^*\leftarrow\argmin_{v\in\mathrm{feasible}}\Delta_{\mathrm{cost}}(v\mid\pi)$
  \Comment{$do_i$ hợp lệ chỉ khi $pu_i\in\pi$}
  \State Thêm $v^*$ vào $\pi$
\EndWhile
\State \Return $\pi$
\end{algorithmic}
\end{algorithm}

%% ================================================================
\section{Bài toán Định tuyến Xe Động (Dynamic VRP)}

\subsection{VRPTW --- ILP với Subtour Elimination tường minh}

\begin{definitionbox}{Vehicle Routing Problem with Time Windows (VRPTW)}
Tìm lộ trình $\{\pi_k\}$ sao cho mỗi khách thăm đúng 1 lần,
trong cửa sổ $[a_i,b_i]$, tổng tải $\leq Q_k$, tổng chi phí nhỏ nhất.
\end{definitionbox}

Hàm mục tiêu:
\begin{equation}
  \min\sum_k\sum_{(i,j)\in\mathcal{E}}c_{ij}x_{ijk}
\end{equation}

Ràng buộc luồng và thời gian:
\begin{align}
  \sum_k\sum_j x_{ijk} &= 1 &&\forall i\in\mathcal{C} \tag{R1}\\
  \sum_i q_i\sum_j x_{ijk} &\leq Q_k &&\forall k \tag{R4}\\
  s_{ik}+\tau_{ij} &\leq s_{jk}+M(1-x_{ijk}) &&\forall i,j,k \tag{R5}\\
  a_i \leq s_{ik} &\leq b_i &&\forall i,k \tag{R6}
\end{align}

\begin{remark}[Subtour Elimination --- MTZ đầy đủ]
Ràng buộc R5 chỉ loại subtour \emph{ngầm} qua time consistency.
Khi cửa sổ thời gian rộng, cần thêm \textbf{Miller--Tucker--Zemlin (MTZ)}:
\begin{align}
  u_{ik} - u_{jk} + n\,x_{ijk} &\leq n-1 &&\forall i\neq j,\;i,j\in\mathcal{C},\;k
  \tag{MTZ}\\
  1 &\leq u_{ik} \leq |\mathcal{C}| &&\forall i\in\mathcal{C},\;k
  \tag{MTZ-bound}\\
  u_{0k} &= 0 &&\forall k
  \tag{MTZ-depot}
\end{align}
$u_{ik}$ là thứ tự (position) node $i$ trên tuyến $k$. Ràng buộc bound và
depot \textbf{bắt buộc} phải có: thiếu chúng, MTZ không loại hết subtour
khi depot không tham gia vào subtour bị loại.
Hoặc dùng \textbf{DFJ cuts} mạnh hơn (nhưng số lượng exponential):
$\sum_{i\in S,j\notin S}x_{ijk}\geq 1$ cho mọi $\emptyset\neq S\subsetneq\mathcal{C}$,
thường được thêm on-the-fly dưới dạng cutting planes trong Branch-and-Cut.
\end{remark}

\subsection{Dynamic VRP --- Công thức MDP}

\begin{definitionbox}{DVRP dưới dạng Markov Decision Process}
$\mathcal{M}=(\mathcal{S},\mathcal{A},\mathcal{P},\mathcal{R},\gamma)$:
\begin{itemize}
  \item $s_t$: vị trí tài xế, đơn chờ, đơn đang giao, SoC đội xe, giao thông.
  \item $a_t$: gán đơn, cập nhật lộ trình, lên lịch sạc.
  \item $r_t = -\mathrm{cost}(a_t) - \lambda\,\mathrm{penalty}(a_t)$.
  \item $\gamma\in[0.9,0.99]$.
\end{itemize}
\end{definitionbox}

\begin{formulabox}{PPO Objective với GAE}
\begin{equation}
  \mathcal{L}^{\mathrm{PPO}}(\theta)
  = \E_t\!\left[
      \min\!\bigl(r_t(\theta)\hat{A}_t,\;
        \mathrm{clip}(r_t(\theta),1{-}\varepsilon,1{+}\varepsilon)\hat{A}_t\bigr)
    \right]
\end{equation}
\begin{equation}
  r_t(\theta)=\frac{\pi_\theta(a_t|s_t)}{\pi_{\theta_\mathrm{old}}(a_t|s_t)},
  \qquad
  \hat{A}_t = \sum_{l=0}^{T-t}(\gamma\lambda_{\mathrm{gae}})^l\delta_{t+l},
  \qquad
  \delta_t = r_t+\gamma V(s_{t+1})-V(s_t)
\end{equation}
$T$ là truncation horizon; $\hat{A}_t$ được chuẩn hóa
$\leftarrow (\hat{A}_t - \mu_A)/\sigma_A$ để giảm variance.
\end{formulabox}

\subsection{ST-GNN cho ETA Prediction}

\begin{formulabox}{Message Passing (ST-GNN) + Quantile Loss}
\begin{align}
  \mathbf{m}_{ij}^{(l)} &= \phi^{(l)}(\mathbf{h}_i^{(l)},\mathbf{h}_j^{(l)},\mathbf{e}_{ij})\\
  \mathbf{h}_i^{(l+1)} &= \psi^{(l)}\!\left(\mathbf{h}_i^{(l)},\bigoplus_{j\in\mathcal{N}(i)}\mathbf{m}_{ij}^{(l)}\right)\\
  \hat{\tau}(u,v,t) &= \mathrm{MLP}\bigl([\mathbf{h}_u^{(L)}\|\mathbf{h}_v^{(L)}\|\mathbf{e}_{uv}\|\mathbf{t}_{\mathrm{emb}}]\bigr)
\end{align}
\begin{equation}
  \mathcal{L}_{\mathrm{ETA}}
  = \mathcal{L}_{\mathrm{Huber}}(\hat{\tau},\tau^*)
  + \eta\sum_{q\in\{0.1,0.5,0.9\}}\mathcal{L}_q(\hat{\tau}_q,\tau^*)
\end{equation}
\end{formulabox}

%% ================================================================
\section{Điều phối Tài xế \& Ghép cặp Tối ưu}

\subsection{Assignment Problem --- Trường hợp Mất cân bằng}

\begin{formulabox}{Hàm chi phí ghép cặp đa tiêu chí}
\begin{equation}
  c_{kj}
  = w_1\hat{\tau}(\ell_k,pu_j)
  + w_2\,\mathrm{urgency}(o_j)
  + w_3\,\mathrm{compat}(d_k,o_j)
  + w_4\,\mathrm{fairness}(d_k),
  \qquad
  \mathrm{urgency}(o_j)=\frac{b_j-t_{\mathrm{now}}}{b_j-a_j}
\end{equation}
\end{formulabox}

\begin{remark}[Trường hợp mất cân bằng và Auction Algorithm]
Thực tế $|\mathcal{D}_{\mathrm{avail}}|\neq|\mathcal{O}_{\mathrm{pending}}|$.
Giải pháp:
\begin{enumerate}[leftmargin=2em]
  \item \textbf{Padding dummy:} thêm node ảo với chi phí lớn để vuông hóa ma trận.
  \item \textbf{Auction Algorithm} (Bertsekas 1988): mỗi đơn ``đấu giá'' tài xế,
        độ phức tạp $O(n^2\log n)$, song song hóa tốt hơn Hungarian.
  \item \textbf{Successive Shortest Path (SSP)} cho min-cost flow: tổng quát nhất.
\end{enumerate}
Với $n=1000$ tài xế, target latency 10ms đòi hỏi \textbf{batched matching}
(chia vùng địa lý) hoặc \textbf{approximate methods} (beam search, greedy).
\end{remark}

\begin{algorithm}[H]
\caption{Hungarian (Kuhn--Munkres) --- $O(n^3)$}
\begin{algorithmic}[1]
\Require $C\in\mathbb{R}^{n\times n}$ (pad dummy nếu cần)
\Ensure $\sigma^*$
\State Row reduce; Col reduce
\Repeat
  \State Tìm phủ tối thiểu $k$ đường
  \If{$k=n$} \textbf{break} \EndIf
  \State $\theta\leftarrow\min$ phần tử không phủ; cập nhật $C$
\Until{phủ hoàn toàn}
\State \Return $\sigma^*$ từ các số 0 độc lập
\end{algorithmic}
\end{algorithm}

\subsection{Dự báo Cầu: ConvLSTM + GAT}

\begin{align}
  \mathbf{H}_t &= \mathrm{ConvLSTM}(\mathbf{H}_{t-1},\mathbf{X}_t)\\
  \mathbf{Z}_t &= \mathrm{GAT}(\mathbf{H}_t,A_{\mathrm{grid}})\\
  \hat{D}_{g,t+\Delta t} &= \mathrm{MLP}(\mathbf{Z}_t[g]\|\mathrm{embed}(t{+}\Delta t))
\end{align}

\subsection{Tái phân bổ tài xế (Repositioning)}

\begin{definitionbox}{Time-Expanded Transportation Problem}
Mở rộng transportation problem sang mạng time-expanded
$G^T=(\mathcal{V}\times\mathcal{T},\mathcal{E}^T)$ để tính thời gian di chuyển
giữa các ô lưới:
\begin{equation}
  \min\sum_{g,g',t}c_{gg'}\,f_{gg'}^t
  \quad\text{s.t.}\quad
  \sum_{g',t'}f_{gg'}^t\leq\hat{S}_g^t,\quad
  \sum_{g,t}f_{gg'}^{t-\tau_{gg'}}\geq\max(0,\hat{D}_{g'}^{t'}-\hat{S}_{g'}^{t'})
\end{equation}
$\tau_{gg'}$ là thời gian di chuyển từ ô $g$ sang $g'$ --- thiếu yếu tố này
trong formulation một giai đoạn là không thực tế.
\end{definitionbox}

\subsection{Đội xe Dị chủng (Heterogeneous Fleet)}
\label{sec:hetfleet}

Thay $Q_k\in\mathbb{R}$ (scalar) bằng vector năng lực:
\begin{equation}
  \mathbf{Q}_k = (\text{seats}_k,\;\text{range}_k,\;\text{cost\_per\_km}_k,\;
                  \text{charge\_speed}_k,\;\text{class}_k)
\end{equation}

\begin{table}[h!]
\centering
\caption{Lớp xe trong đội Green SM (ví dụ)}
\renewcommand{\arraystretch}{1.2}
\begin{tabular}{lcccc}
\toprule
\textbf{Xe} & \textbf{Chỗ ngồi} & \textbf{Range (km)} & \textbf{Sạc nhanh} & \textbf{Lớp dịch vụ}\\
\midrule
VinFast Bike & 1 & 80 & Không & Bike/Delivery\\
VF5           & 4 & 320 & Có & Standard\\
VF e34        & 5 & 380 & Có & Standard/Premium\\
VF8           & 7 & 420 & Có (150 kW) & Premium\\
\bottomrule
\end{tabular}
\end{table}

Ma trận tương thích giữa loại xe $k$ và loại yêu cầu $j$:
\begin{equation}
  \mathrm{compat\_fleet}(k,j) =
  \begin{cases}
    1 & \text{nếu }\mathrm{class}_k \geq \mathrm{required\_class}(j)
        \;\wedge\; \mathrm{seats}_k \geq q_j\\
    0 & \text{ngược lại}
  \end{cases}
\end{equation}

\subsection{Driver Fairness --- Hình thức hóa}

\begin{formulabox}{Các thước đo công bằng cho tài xế}
\textbf{(a) Hệ số Gini thu nhập trong ca:}
\begin{equation}
  \Gini = \frac{\sum_{k,k'}|R_k - R_{k'}|}{2m\sum_k R_k},
  \quad R_k = \text{thu nhập tài xế } k \text{ trong ca}
\end{equation}
\textbf{(b) Lexicographic Max-Min:} tối ưu lần lượt theo thứ tự từ tài xế
thu nhập thấp nhất:
$\max\min_k R_k$, subject to các ràng buộc routing.\\
\textbf{(c) Nash Bargaining Solution:}
$\max\prod_k(R_k - R_k^0)$ với $R_k^0$ là thu nhập tham chiếu (outside option).
\end{formulabox}

Hàm fairness penalty tích hợp vào cost matrix:
$\mathrm{fairness}(d_k) = \max(0,\,\bar{R} - R_k^{\mathrm{current}})$,
ưu tiên tài xế đang có thu nhập thấp hơn trung bình.

\subsection{Ràng buộc Thực tế Nâng cao}

\subsubsection{Soft Constraints + Augmented Lagrangian}
\begin{equation}
  \mathcal{L}_{\mathrm{total}}
  = \mathcal{L}_{\mathrm{cost}}+\sum_c\lambda_c\,\mathrm{violation}_c(\bm{x}),
  \qquad
  \lambda_c^{(t+1)}=\lambda_c^{(t)}+\rho\cdot\mathrm{violation}_c(\bm{x}^{(t)})
\end{equation}

\subsubsection{Robust Optimization --- DRO \& CVaR}

\begin{remark}[Phân phối thực tế của thời gian di chuyển]
Giả định $\tau\sim\mathcal{N}(\hat{\tau},\sigma^2)$ không phù hợp thực tế:
phân phối thời gian di chuyển đô thị Việt Nam là \textbf{heavy-tailed, right-skewed}
(gần Lognormal hoặc Gamma). Dùng hai phương pháp thay thế:
\end{remark}

\textbf{(a) Distributionally Robust Optimization (DRO)} với Wasserstein ball:
\begin{equation}
  \min_{\bm{x}}\max_{\mathbb{Q}\in\mathcal{B}_\varepsilon(\hat{\mathbb{P}})}
  \E_{\mathbb{Q}}[\ell(\bm{x},\xi)]
\end{equation}
$\mathcal{B}_\varepsilon(\hat{\mathbb{P}})=\{\mathbb{Q}:\mathcal{W}_p(\mathbb{Q},\hat{\mathbb{P}})\leq\varepsilon\}$
là Wasserstein ball bán kính $\varepsilon$ quanh phân phối thực nghiệm $\hat{\mathbb{P}}$
(Mohajerin Esfahani \& Kuhn 2018).

\textbf{(b) CVaR Chance Constraint} (Rockafellar \& Uryasev 2000):
\begin{equation}
  \CVaR_\alpha[\tau(o_i)] \leq b_i - t_{\mathrm{now}}
\end{equation}
\begin{equation}
  \CVaR_\alpha[\tau] = \min_{v}\left\{v + \frac{1}{1-\alpha}\E[\max(\tau-v,0)]\right\}
\end{equation}
$\alpha\in[0.9,0.99]$ kiểm soát mức độ bảo thủ của lịch giao hàng.

\subsection{Lịch làm việc Tài xế \& Giới hạn Mệt mỏi (Crew Scheduling)}

Mỗi tài xế Green SM / Xanh SM chịu ràng buộc pháp lý và an toàn:
tổng thời gian lái tối đa $H_{\max}=8$ giờ/ngày, nghỉ bắt buộc sau mỗi
$\Delta^{\mathrm{drive}}=4$ giờ liên tục.

\begin{definitionbox}{Bài toán Crew Scheduling (CS)}
\textbf{Biến quyết định:}
\begin{itemize}
  \item $z_{k\sigma}\in\{0,1\}$: tài xế $k$ được gán vào ca làm việc (shift) $\sigma$.
  \item $t_k^{\mathrm{start}}, t_k^{\mathrm{end}}$: thời điểm bắt đầu/kết thúc ca.
  \item $\delta_k^{(m)}$: khoảng thời gian lái liên tục thứ $m$ của tài xế $k$.
\end{itemize}

\textbf{Mục tiêu:} Tối thiểu tổng chi phí nhân sự và thời gian không có tài xế:
\begin{equation}
  \min \sum_k \sum_\sigma c_\sigma z_{k\sigma}
  + \mu \sum_{t}\max\!\left(0,\, D(t) - \sum_k A_k(t)\right)
\end{equation}
$D(t)$ = nhu cầu tài xế dự báo tại thời điểm $t$;
$A_k(t)=1$ nếu tài xế $k$ đang trong ca và chưa vượt giới hạn mệt mỏi.
\end{definitionbox}

\begin{formulabox}{Ràng buộc Mệt mỏi \& Thời gian lái}
\begin{align}
  \sum_\sigma z_{k\sigma} &\leq 1 &&\forall k \tag{CS1: Mỗi tài xế 1 ca/ngày}\\
  t_k^{\mathrm{end}} - t_k^{\mathrm{start}} &\leq H_{\max} &&\forall k \tag{CS2: Tổng giờ lái}\\
  \sum_m \delta_k^{(m)} &\leq H_{\max} &&\forall k \tag{CS3: Tổng tích lũy}\\
  \delta_k^{(m)} &\leq \Delta^{\mathrm{drive}} &&\forall k,m \tag{CS4: Giới hạn liên tục}
\end{align}
\end{formulabox}

\textbf{Phương pháp giải:}

\begin{itemize}
  \item \textbf{Set Partitioning / Column Generation:} Sinh ca khả dĩ (columns) qua
    subproblem shortest-path trên đồ thị thời gian; giải master LP bằng simplex.
  \item \textbf{Time-Space Network:} Mô hình hóa trạng thái tài xế như nút
    $(\ell_k, t, \delta_k^{\mathrm{accum}})$ trên đồ thị không-thời gian; tìm đường
    đi chi phí tối thiểu.
  \item \textbf{RL-based Fatigue Model:} Mô hình hóa mức độ mệt mỏi như trạng thái
    ẩn trong Markov Decision Process, điều chỉnh hành vi đề xuất ca theo trạng thái
    thực tế (tốc độ phản xạ, tỷ lệ dừng đột ngột từ dữ liệu cảm biến xe).
\end{itemize}

%% ================================================================
\section{Tối ưu hóa Hệ thống Xe Điện (EV-Specific)}

\subsection{Mô hình Tiêu thụ Năng lượng}

\begin{evbox}{Energy Consumption Model (phi tuyến)}
\begin{equation}
  h_{ij} = \alpha\,d_{ij} + \beta\,\Delta h_{ij} + \gamma\,v_{ij}^2\,d_{ij}
           + \delta\,m_{\mathrm{load}}\,d_{ij} + \epsilon_{\mathrm{HVAC}}
\end{equation}
\begin{itemize}
  \item $\alpha d_{ij}$: tiêu thụ cơ bản theo quãng đường (kWh/km)
  \item $\beta\Delta h_{ij}$: năng lượng leo dốc (có thể âm khi xuống dốc)
  \item $\gamma v_{ij}^2 d_{ij}$: lực cản khí động học (tốc độ cao tốn điện nhiều)
  \item $\delta m_{\mathrm{load}} d_{ij}$: tải hàng hóa / hành khách
  \item $\epsilon_{\mathrm{HVAC}}$: điều hòa nhiệt độ (phụ thuộc thời tiết)
\end{itemize}
\end{evbox}

Tham số $(\alpha,\beta,\gamma,\delta)$ hiệu chỉnh (calibrate) trên dữ liệu telemetry
thực tế của đội xe VinFast. Mô hình này thay $c_{ij}$ trong VRPTW bằng
chi phí năng lượng $h_{ij}$.

\subsection{E-VRP-CS --- Electric VRP with Charging Stations}

\begin{evbox}{E-VRP-CS: Mở rộng VRPTW với SoC và trạm sạc}
Mở rộng tập node: $\mathcal{V}' = \{0\}\cup\mathcal{C}\cup\mathcal{F}$
($\mathcal{F}$ = tập trạm sạc V-Green).
Biến mới:
$e_{ki}\in[E_k^{\min},E_k^{\max}]$ = SoC xe $k$ khi đến node $i$;
$y_{kf}\geq 0$ = lượng điện nạp tại trạm $f$.
\end{evbox}

\subsubsection{Ràng buộc SoC}

Đặt biến phụ $p_{ijk} = h_{ij}\cdot x_{ijk}$ (tuyến hóa tích nhị phân--liên tục)
với các ràng buộc McCormick:
\begin{align}
  p_{ijk} &\leq h_{ij} &&\forall i,j,k \tag{P1}\\
  p_{ijk} &\leq M\cdot x_{ijk} &&\forall i,j,k \tag{P2}\\
  p_{ijk} &\geq h_{ij} - M(1-x_{ijk}) &&\forall i,j,k \tag{P3}\\
  p_{ijk} &\geq 0 &&\forall i,j,k \tag{P4}
\end{align}

Dùng $p_{ijk}$ thay thế, các ràng buộc SoC trở thành tuyến tính:
\begin{align}
  e_{kj} &\leq e_{ki} - p_{ijk} + M(1-x_{ijk}) &&\forall i,j,k
            \tag{E1a: SoC upper}\\
  e_{kj} &\geq e_{ki} - p_{ijk} - M(1-x_{ijk}) &&\forall i,j,k
            \tag{E1b: SoC lower}\\
  E_k^{\min} &\leq e_{ki} \leq E_k^{\max} &&\forall i,k
              \tag{E2: SoC bound}\\
  e_{kf}^{\mathrm{after}} &= e_{kf}^{\mathrm{before}} + y_{kf} &&\forall f,k
                            \tag{E3: Charging gain}\\
  y_{kf} &\leq E_k^{\max} - e_{kf}^{\mathrm{before}} &&\forall f,k
          \tag{E4: Không sạc quá đầy}
\end{align}

\begin{remark}
Khi $x_{ijk}=1$: E1a và E1b kẹp $e_{kj}=e_{ki}-h_{ij}$ (tiêu thụ đúng).
Khi $x_{ijk}=0$: cả hai ràng buộc tự động relaxed bởi $\pm M$, không ràng buộc $e_{kj}$.
Điều này đúng về mặt logic và khác hoàn toàn với cách viết $e_{kj}=\ldots+M(1-x_{ijk})$
(dấu $=$ kết hợp big-M là không hợp lệ vì big-M chỉ làm relaxation một chiều).
\end{remark}

\subsubsection{Hàm mục tiêu E-VRP-CS}

\begin{equation}
  \min\underbrace{\sum_{k,i,j}c_{ij}x_{ijk}}_{\text{Vận chuyển}}
  + \underbrace{\sum_{k,f}\rho_f^{\mathrm{elec}}\,y_{kf}}_{\text{Chi phí điện}}
  + \underbrace{\sum_{k,f}\beta_{\mathrm{deg}}\,y_{kf}}_{\text{Hao mòn pin}}
  + \lambda\underbrace{\sum_i\max(0,t_i-b_i)}_{\text{Trễ hạn}}
\end{equation}

$\rho_f^{\mathrm{elec}}$ phụ thuộc giờ sạc (off-peak rẻ hơn peak),
$\beta_{\mathrm{deg}}$ là chi phí hao mòn pin mỗi kWh nạp.

\subsection{Đường cong Sạc Phi tuyến (Piecewise Linear)}

Tốc độ sạc chậm dần khi SoC $> 80\%$ (CC-CV protocol):
\begin{equation}
  \frac{de}{dt} =
  \begin{cases}
    P_{\max}/E_{\max} & e < 0.8\,E_{\max} \quad\text{(CC phase)}\\
    P_{\max}/E_{\max}\cdot(1 - e/E_{\max})^2 & e \geq 0.8\,E_{\max} \quad\text{(CV phase)}
  \end{cases}
\end{equation}

Xấp xỉ tuyến tính từng đoạn (piecewise linear) với $K_c$ breakpoints
$0=e_0<e_1<\cdots<e_{K_c}=E_{\max}$ và độ dốc $\mu_s$ trên đoạn $[e_{s-1},e_s]$:

\begin{formulabox}{Piecewise Linear Charging với SOS2}
\begin{align}
  t_{\mathrm{charge}}(y) &= \sum_{s=1}^{K_c}\mu_s\,\lambda_s \label{eq:pwl-time}\\
  y &= \sum_{s=1}^{K_c}\lambda_s\,\Delta e_s, \quad \Delta e_s = e_s - e_{s-1}\\
  \sum_{s=1}^{K_c}\lambda_s &= 1, \quad \lambda_s \geq 0 \quad\forall s
  \tag{Convex combination}\\
  \{\lambda_s\} &\in \mathrm{SOS2} \tag{SOS2: chỉ 2 biến liền kề $\neq 0$}
\end{align}
\end{formulabox}

Ràng buộc SOS2 (Special Ordered Set type 2) đảm bảo chỉ tối đa 2 biến
$\lambda_s, \lambda_{s+1}$ liền kề có thể khác 0 cùng lúc, giữ cho
xấp xỉ luôn nằm trên đúng một đoạn tuyến tính của đường cong CC-CV.
Không có SOS2, bộ giải có thể chọn tổ hợp $\lambda$ không liền kề
cho ra nghiệm vô nghĩa về vật lý.

\subsection{Charging Station Queuing \& Reservation}

\begin{evbox}{Mô hình hàng đợi M/M/c tại trạm sạc}
Trạm sạc $f$ có $c_f$ trụ. Xe đến theo Poisson($\lambda_f$),
thời gian sạc $\sim\mathrm{Exp}(\mu_f)$. Mô hình M/M/c:
\begin{equation}
  \Prob(\text{đợi}) = \frac{(c_f\rho_f)^{c_f}}{c_f!\,(1-\rho_f)}\cdot
  \left[\sum_{n=0}^{c_f-1}\frac{(c_f\rho_f)^n}{n!}
        + \frac{(c_f\rho_f)^{c_f}}{c_f!\,(1-\rho_f)}\right]^{-1}
  \quad (\text{Erlang C})
\end{equation}
$\rho_f = \lambda_f/(c_f\mu_f) < 1$. Thời gian chờ trung bình:
$W_q = \Prob(\text{đợi})/(\mu_f(c_f-\lambda_f/\mu_f))$.
\end{evbox}

\subsubsection{Charging Slot Reservation}

\begin{equation}
  \min\sum_{k,f,t}w_{kft}\,\mathrm{wait}_{kft}
  \quad\text{s.t.}\quad
  \sum_k z_{kft} \leq c_f \;\forall f,t
  \quad\text{(Capacity trụ sạc)}
\end{equation}

$z_{kft}\in\{0,1\}$: xe $k$ chiếm trụ tại trạm $f$ trong slot $t$.
Hệ thống đặt chỗ trước (reservation) giảm $W_q$ xuống gần 0 cho các xe có kế hoạch.

\subsection{Battery Degradation \& Lifecycle Cost}

\begin{equation}
  C_{\mathrm{deg}} = \underbrace{c_{\mathrm{bat}}\cdot\frac{1}{L_{\mathrm{cycle}}}}_{\mathrm{Rainflow\ counting}}
  \quad\text{với}\quad
  L_{\mathrm{cycle}}(DoD) = A\cdot e^{-B\cdot DoD}
\end{equation}

$DoD$ = depth of discharge mỗi chu kỳ sạc, $c_{\mathrm{bat}}$ = chi phí thay pin,
$L_{\mathrm{cycle}}$ = số chu kỳ theo DoD (mô hình Arrhenius). Đưa $C_{\mathrm{deg}}$
vào objective của E-VRP-CS khuyến khích sạc nhanh-nông thay vì sạc sâu.

\subsection{Risk-Averse RL cho Range Anxiety}

\emph{Range anxiety}: tâm lý sợ hết pin trước khi đến trạm sạc.
Giải quyết bằng \textbf{CVaR-PPO}:

\begin{formulabox}{CVaR-PPO cho an toàn năng lượng}
\begin{equation}
  \mathcal{L}^{\mathrm{CVaR-PPO}}(\theta)
  = \mathcal{L}^{\mathrm{PPO}}(\theta)
  - \kappa\,\CVaR_\alpha\!\left[\min_{t\leq T}e_k(t)\right]
\end{equation}
\begin{equation}
  \text{Reward shaping: }\quad
  r_t \leftarrow r_t - \eta\,\mathbb{1}[e_k(t) < E_k^{\mathrm{safe}}]
  \cdot \left(\frac{E_k^{\mathrm{safe}} - e_k(t)}{E_k^{\mathrm{max}}}\right)^2
\end{equation}
\end{formulabox}

$E_k^{\mathrm{safe}}$ là ngưỡng an toàn (ví dụ 15\% SoC). CVaR kiểm soát
phân vị $\alpha$ của SoC tối thiểu, tránh chính sách ``liều'' cạn pin
dù kỳ vọng phần thưởng cao.

\subsection{Quy hoạch Mạng lưới Trạm Sạc (Charging Network Planning)}

V-Green cần xác định \emph{vị trí} và \emph{công suất} cho mạng lưới trạm sạc mới
trong các khu đô thị Vinhomes để tối thiểu chi phí đầu tư trong khi đảm bảo
mức độ phục vụ (coverage) cho toàn đội xe.

\begin{definitionbox}{Bài toán Facility Location cho Trạm Sạc}
\textbf{Tập hợp:}
\begin{itemize}
  \item $\mathcal{F}$: tập vị trí ứng viên (candidate sites).
  \item $\mathcal{D}$: tập điểm nhu cầu (demand points, ví dụ: khu dân cư, tòa nhà văn phòng).
  \item $K$: tổng ngân sách trạm sạc.
\end{itemize}

\textbf{Biến quyết định:}
$y_f\in\{0,1\}$: có đặt trạm tại vị trí $f$ không;
$n_f\in\mathbb{Z}^+$: số đầu sạc tại $f$;
$z_{df}\in[0,1\}$: tỷ lệ nhu cầu tại $d$ được phục vụ bởi $f$.

\textbf{Mô hình:}
\begin{align}
  \min\quad & \sum_{f\in\mathcal{F}} c_f^{\mathrm{fixed}} y_f
              + \sum_{f\in\mathcal{F}} c_f^{\mathrm{cap}} n_f
              + \sum_{d,f} c_{df}^{\mathrm{travel}} z_{df} \notag\\[4pt]
  \text{s.t.}\quad
  & \sum_{f\in\mathcal{F}} z_{df} \geq 1 - \epsilon_d \quad\forall d \tag{FL1: Coverage}\\
  & z_{df} \leq y_f \quad\forall d,f \tag{FL2: Open facility}\\
  & \sum_d \lambda_d z_{df} \leq \mu_f n_f \quad\forall f \tag{FL3: Capacity}\\
  & \sum_{f} y_f \leq K \tag{FL4: Budget}\\
  & \tau_{df} \leq R_{\max} \;\Rightarrow\; z_{df}\text{ feasible} \tag{FL5: Range}
\end{align}
\textbf{FL5} loại bỏ kết nối $d\to f$ nếu khoảng cách $\tau_{df}$ vượt tầm hoạt động
của xe $R_{\max}$.
\end{definitionbox}

\textbf{Phương pháp giải:}
\begin{itemize}
  \item \textbf{Capacitated Facility Location (CFL)}: NP-hard; giải bằng Branch \& Bound
    hoặc Lagrangian Relaxation.
  \item \textbf{$p$-Median Problem}: chọn đúng $p=K$ vị trí tối thiểu tổng khoảng cách
    có trọng số; giải bằng heuristic Local Search hoặc GRASP.
  \item \textbf{Simulation-Optimization}: Kết hợp Monte Carlo simulation (mô phỏng
    các chuyến đi EV theo phân phối cầu thực tế) với metaheuristic (SA, GA) để
    đánh giá phương án vị trí.
\end{itemize}

\subsection{Bài toán Vị trí Trạm Đổi Pin (Battery Swap Station Location)}

Đổi pin (battery swap) như mô hình VinFast / NIO giải quyết điểm yếu về
thời gian sạc lâu: thay vì sạc 30--60 phút, tài xế đổi pin trong $<5$ phút.

\begin{definitionbox}{Battery Swap Station Location Problem (BSSLP)}
\textbf{Đặc thù bổ sung so với Charging Network Planning:}
\begin{itemize}
  \item Mỗi trạm phải duy trì một \emph{kho pin} (battery inventory)
    $I_f(t)$: số pin đầy sẵn sàng tại trạm $f$, thời điểm $t$.
  \item Ràng buộc cân bằng kho:
    \begin{equation}
      I_f(t+1) = I_f(t) + R_f(t) - S_f(t)
    \end{equation}
    $R_f(t)$: pin được sạc xong (replenishment); $S_f(t)$: pin xuất cho xe.
  \item \textbf{Stockout penalty:} $P_{\mathrm{so}}$ nếu $I_f(t)=0$ và có xe
    đến đổi -- tái hiện bài toán newsvendor.
\end{itemize}

\textbf{Mô hình tích hợp Vị trí + Tồn kho:}
\begin{equation}
  \min\;\sum_f \left[c_f^{\mathrm{loc}}\,y_f
    + c^{\mathrm{bat}}\,B_f
    + \E_\xi\!\left(P_{\mathrm{so}}\,\mathrm{unmet}_f(\xi)\right)\right]
\end{equation}
$B_f$ = số pin dự trữ ban đầu tại $f$; $\xi$ = vector nhu cầu ngẫu nhiên.
\end{definitionbox}

\textbf{Newsvendor cho mỗi trạm:}
\begin{equation}
  B_f^* = F^{-1}_f\!\left(\frac{c^{\mathrm{sold}} - c^{\mathrm{bat}}}{c^{\mathrm{sold}} + P_{\mathrm{so}}}\right)
\end{equation}
$F_f$: CDF của nhu cầu đổi pin tại trạm $f$.

\subsection{Vehicle-to-Grid (V2G) --- Mô hình hóa}

Xe điện không chỉ là phương tiện tiêu thụ điện -- khi đỗ, chúng có thể
\emph{phát điện ngược lưới} (V2G, Vehicle-to-Grid), tạo nguồn thu bổ sung
và hỗ trợ cân bằng lưới điện quốc gia.

\begin{evbox}{Mô hình tối ưu V2G}
Mỗi xe $k$ trong khoảng thời gian đỗ $[t^a_k, t^d_k]$ (arrive / depart) có thể:
\begin{itemize}
  \item \textbf{Nạp điện (G2V):} $p_k^{\mathrm{G2V}}(t) \geq 0$
  \item \textbf{Phát điện ngược (V2G):} $p_k^{\mathrm{V2G}}(t) \geq 0$
\end{itemize}

\textbf{SoC dynamics với V2G:}
\begin{equation}
  e_k(t+1) = e_k(t) + \eta^+\, p_k^{\mathrm{G2V}}(t)\,\Delta t
              - \frac{1}{\eta^-}\,p_k^{\mathrm{V2G}}(t)\,\Delta t
\end{equation}
$\eta^+, \eta^-$: hiệu suất nạp/xả.

\textbf{Ràng buộc:}
\begin{align}
  e_k^{\min} &\leq e_k(t) \leq e_k^{\max} &&\forall k,t \tag{V1: SoC bounds}\\
  p_k^{\mathrm{G2V}}(t) + p_k^{\mathrm{V2G}}(t) &\leq P_k^{\max} &&\forall k,t \tag{V2: Power limit}\\
  e_k(t^d_k) &\geq e_k^{\mathrm{depart}} &&\forall k \tag{V3: Departure SoC}\\
  p_k^{\mathrm{G2V}}(t)\cdot p_k^{\mathrm{V2G}}(t) &= 0 &&\forall k,t \tag{V4: No simultaneous}
\end{align}
V4 là ràng buộc phi tuyến; linearize bằng biến nhị phân $b_k(t)\in\{0,1\}$:
\begin{equation}
  p_k^{\mathrm{G2V}}(t) \leq P_k^{\max} b_k(t),\quad
  p_k^{\mathrm{V2G}}(t) \leq P_k^{\max}(1-b_k(t))
\end{equation}

\textbf{Mục tiêu tối ưu hóa toàn hệ thống:}
\begin{equation}
  \max\;\sum_k\sum_t\left[
    \lambda_t^{\mathrm{sell}}\,p_k^{\mathrm{V2G}}(t)
    - \lambda_t^{\mathrm{buy}}\,p_k^{\mathrm{G2V}}(t)
  \right]\Delta t
\end{equation}
$\lambda_t^{\mathrm{sell}}, \lambda_t^{\mathrm{buy}}$: giá điện bán/mua tại thời điểm $t$ (time-of-use tariff).
\end{evbox}

\textbf{Ứng dụng thực tế:} Kết hợp lịch sạc M/M/c (§5.4) với V2G, hệ thống
có thể tham gia thị trường điều tiết tần số (frequency regulation market),
tạo nguồn doanh thu phụ cho V-Green đồng thời hỗ trợ lưới điện quốc gia
trong giờ cao điểm.

%% ================================================================
\section{Tối ưu hóa Ride-hailing Chở Khách}

\subsection{So sánh Delivery vs Ride-hailing}

\begin{table}[h!]
\centering
\caption{Phân biệt Delivery (logistics) và Ride-hailing (hành khách)}
\renewcommand{\arraystretch}{1.3}
\begin{tabular}{lp{4.5cm}p{4.5cm}}
\toprule
\textbf{Đặc trưng} & \textbf{Delivery (đã có)} & \textbf{Ride-hailing (mới)} \\
\midrule
Batching & OBP, ST-DBSCAN & Ride-pooling, Shareability Network \\
Định tuyến & VRPTW, DVRP & DARP, RTV-graph \\
Time window & 20--30 phút & 3--10 phút (khách hủy nhanh) \\
Capacity & Thể tích/khối lượng & Số ghế $+$ comfort class \\
Tỷ lệ hủy & Thấp ($< 5\%$) & Cao (8--15\% ở VN) \\
Pricing & Cố định / zone & Surge real-time \\
Detour & Không relevant & Ràng buộc cứng/mềm \\
Thực thể chờ & Đơn hàng & Hành khách (có tâm lý) \\
\bottomrule
\end{tabular}
\end{table}

\subsection{Dial-a-Ride Problem (DARP)}

\begin{definitionbox}{DARP --- Bài toán nền tảng cho Ride-hailing}
DARP (Cordeau \& Laporte 2007): mở rộng VRPTW với ràng buộc:
\begin{itemize}
  \item \textbf{Precedence:} $pickup_i$ trước $dropoff_i$ cho từng khách $i$.
  \item \textbf{Max ride time:} thời gian hành khách trên xe bị giới hạn:
        $s_{i^-,k} - s_{i^+,k} \leq L_i$ ($i^+$ = pickup, $i^-$ = dropoff).
  \item \textbf{Pairing:} pickup và dropoff cùng xe.
\end{itemize}
\end{definitionbox}

\subsubsection{Công thức ILP chuẩn (Cordeau 2006)}

Hàm mục tiêu:
\begin{equation}
  \min\sum_k\left(c_0 T_k + \sum_{i\in\mathcal{C}}d_i w_i\right)
\end{equation}
$T_k$ = tổng thời gian tuyến $k$; $w_i$ = thời gian chờ của khách $i$;
$c_0$ = chi phí vận hành, $d_i$ = trọng số khó chịu của khách.

Ràng buộc bổ sung so với VRPTW:

\medskip
\noindent\textbf{Ký hiệu phụ:} Đặt $\mathcal{C}^+_k(t) \subseteq \mathcal{C}$ là tập hành khách đang ở trên xe $k$ tại thời điểm $t$, tức là những khách đã được đón (pickup) nhưng chưa được trả (dropoff):
\[
  \mathcal{C}^+_k(t) \;=\; \bigl\{\,i\in\mathcal{C} : s_{i^+,k} \leq t < s_{i^-,k}\bigr\}.
\]

\begin{align}
  s_{i^-,k} - s_{i^+,k} &\leq L_i &&\forall i,k \tag{D1: Max ride time}\\
  s_{i^+,k} &\leq b_{i^+} &&\forall i,k \tag{D2: Pickup window}\\
  \sum_k\left(\sum_j x_{i^+jk} - \sum_j x_{ji^-k}\right) &= 0 &&\forall i \tag{D3: Pairing}\\
  \sum_{i\in\mathcal{C}^+_k(t)}q_i &\leq Q_k &&\forall k,\,t \tag{D4: Capacity (seats)}
\end{align}

\subsection{Shareability Network \& Ride-Pooling}

\begin{ridebox}{Shareability Network (Santi et al.~2014, PNAS)}
\textbf{Cấu trúc:} Đồ thị $G_S=(\mathcal{T},\mathcal{E}_S)$ trong đó:
\begin{itemize}
  \item Node $t_i\in\mathcal{T}$: một chuyến đi (trip request).
  \item Edge $(t_i,t_j)\in\mathcal{E}_S$: hai chuyến có thể ghép chung xe
        mà không vi phạm ngưỡng detour của cả hai khách.
\end{itemize}
\textbf{Bài toán:} Maximum Weighted Matching trên $G_S$:
\begin{equation}
  \max\sum_{(t_i,t_j)\in M}w_{ij},
  \quad\text{s.t. mỗi trip thuộc tối đa 1 cặp ghép}
\end{equation}
$w_{ij}$ = tiết kiệm quãng đường khi ghép $t_i$ và $t_j$.
\end{ridebox}

\subsubsection{RTV-Graph: Ghép 2--4 hành khách}

Alonso-Mora et al.\ (2017, PNAS) mở rộng sang ghép $n_{\max}=4$ khách:

\begin{equation}
  G_{\mathrm{RTV}} = (V_R\cup V_T\cup V_V,\;\mathcal{E}_{RT}\cup\mathcal{E}_{TV})
\end{equation}

\begin{itemize}
  \item $V_R$ = requests, $V_T$ = feasible trips (tập hợp requests ghép được),
        $V_V$ = vehicles.
  \item ILP: $\min\sum c_t\,x_t + c_{\mathrm{ignore}}\sum_r(1-\sum_{t\ni r}x_t)$.
  \item Giải gần đúng bằng Branch-and-Price hoặc column generation.
\end{itemize}

\subsection{Detour Constraint \& Passenger Acceptance}

Mỗi khách $i$ chấp nhận đi vòng tối đa $\Delta_i^{\max}$:
\begin{equation}
  \tau_{\mathrm{actual}}(i) - \tau_{\mathrm{direct}}(i) \leq \Delta_i^{\max}
\end{equation}

Mô hình chấp nhận xác suất (discrete choice):
\begin{equation}
  P_i(\mathrm{accept}\mid\text{detour},\text{discount})
  = \frac{1}{1+e^{-(\beta_0 - \beta_1\cdot\mathrm{detour}_i + \beta_2\cdot\mathrm{discount}_i)}}
\end{equation}

Chính sách giá tối ưu: giảm giá theo \% detour để duy trì tỷ lệ chấp nhận:
$\mathrm{discount}_i^* = \argmax_d\left[P_i(\mathrm{accept})\cdot\mathrm{Revenue}(d)\right]$.

\subsection{Cancellation \& No-show Prediction}

\begin{ridebox}{Adjusted Cost Matrix với Cancellation Risk}
Xác suất hủy $\hat{p}_{kj}^{\mathrm{cancel}} = \sigma(\mathbf{f}_{kj}^\top\bm\theta)$ với
features: thời gian chờ ước tính, khoảng cách tài xế, lịch sử khách, thời tiết, giờ.
\begin{equation}
  c_{kj}^{\mathrm{adj}}
  = (1-\hat{p}_{kj})\cdot c_{kj}
  + \hat{p}_{kj}\cdot C_{\mathrm{wasted}}
\end{equation}
$C_{\mathrm{wasted}}$ = chi phí tài xế chạy không về điểm mới (deadhead cost).
\end{ridebox}

\subsubsection{Driver Acceptance Prediction}

Tài xế cũng có quyền từ chối $\Rightarrow$ \textbf{Stochastic Bipartite Matching}:

\begin{equation}
  \E[\text{matched}] = \sum_{(k,j)\in M} P(\text{accept}_{kj})\cdot P(\text{not cancel}_{kj})
\end{equation}

Giải bằng \textbf{Online Bipartite Matching} với Robbins--Monro updates khi xác suất
được học online từ phản hồi thực tế.

\subsection{Dynamic Surge Pricing}

\begin{ridebox}{Spatio-Temporal Surge Pricing Model}
Hệ số surge $\rho_{g,t}\geq 1$ cho ô lưới $g$ tại thời điểm $t$:
\begin{equation}
  \rho_{g,t} = 1 + \kappa\cdot\max\!\left(0,\;\frac{\hat{D}_{g,t}-\hat{S}_{g,t}}{\hat{S}_{g,t}}\right)
\end{equation}
\end{ridebox}

\subsubsection{Mô hình Cân bằng Cung-Cầu}

Độ co giãn cầu theo giá:
$\epsilon_D = \partial Q/\partial P\cdot P/Q < 0$

Độ co giãn cung (tài xế online tăng khi giá tăng):
$\epsilon_S = \partial S/\partial P\cdot P/S > 0$

Giá cân bằng (market clearing):
\begin{equation}
  P^* = \argmin_P\left|D(P) - S(P)\right|
  \quad\Leftrightarrow\quad
  D_0 P^{\epsilon_D} = S_0 P^{\epsilon_S}
  \quad\Rightarrow\quad
  P^* = \left(\frac{D_0}{S_0}\right)^{\frac{1}{\epsilon_S-\epsilon_D}}
\end{equation}

\subsubsection{A/B Testing \& Causal Inference}

Đánh giá tác động thực của surge pricing bằng:
\begin{itemize}
  \item \textbf{Switchback Experiments:} xen kẽ treatment/control theo block thời gian.
  \item \textbf{Synthetic Control:} xây dựng counterfactual từ các ô lưới tương đồng.
  \item Phương trình suy luận nhân quả (DiD):
        $\tau_{\mathrm{ATT}} = \E[Y_{t,g}^{(1)}-Y_{t,g}^{(0)}\mid T_{t,g}=1]$.
\end{itemize}

\subsection{Airport / Hub Queueing \& Last-100m}

\subsubsection{FIFO Queue tại Sân bay}

Tại Nội Bài / Tân Sơn Nhất, cân bằng FIFO fairness vs proximity:

\begin{equation}
  \text{score}(k,j) = w_1\cdot\hat{\tau}(\ell_k,pu_j)
  + w_2\cdot\left(t_{\mathrm{now}} - t_k^{\mathrm{queue}}\right)^{-1}
\end{equation}

Phần tử thứ hai ưu tiên tài xế chờ lâu hơn.

\subsubsection{Last-100m Pickup Matching}

Khu Vinhomes, sân bay, TTTM: nhiều cổng/terminal. Mismatch tốn 30s--2 phút/chuyến.

\begin{equation}
  pu_j^* = \argmin_{p\in\mathcal{P}_j}\left[
    d(\ell_k, p) + \beta\cdot d(\ell_j^{\mathrm{actual}}, p)
  \right]
\end{equation}

$\mathcal{P}_j$ = tập điểm đón khả dĩ (POI cluster), $\ell_j^{\mathrm{actual}}$
= vị trí GPS thực của khách. POI clustering dùng DBSCAN trên lịch sử pickup
kết hợp embedding từ bản đồ vector.

\subsection{Hành khách Đa hạng \& Tối ưu Lexicographic (VIP Tiers)}

Xanh SM vận hành nhiều hạng dịch vụ: Standard (xe thường), Plus (xe rộng),
Premium (VinFast VF8/9, ghế da, wifi). Hệ thống cần đảm bảo khách Premium
được phục vụ \emph{trước} và với xe phù hợp.

\begin{ridebox}{Mô hình Lexicographic Multi-Objective Optimization}
Đặt $\mathcal{K}^{(r)}$ là tập tài xế loại $r$ (Standard $r=1$, Plus $r=2$,
Premium $r=3$). Yêu cầu đặt xe hạng $r$ \emph{chỉ} được ghép với tài xế
loại $r'\geq r$.

\textbf{Phân cấp mục tiêu Lexicographic:}
\begin{equation}
  \mathrm{lex\text{-}min}\;\bigl(W^{(3)},\, W^{(2)},\, W^{(1)}\bigr)
\end{equation}
$W^{(r)}$ = tổng thời gian chờ của khách hạng $r$. Tối ưu hóa lần lượt từ
hạng cao xuống hạng thấp: trước tiên tối thiểu $W^{(3)}$ (Premium), rồi
$W^{(2)}$ trong điều kiện $W^{(3)}$ đã tối ưu, v.v.

\textbf{Ràng buộc hạng xe:}
\begin{align}
  x_{ijk} &= 0 \quad\forall k\in\mathcal{K}^{(r')}: r' < r_i
    \tag{T1: Tier compatibility}\\
  \sum_{k\in\mathcal{K}^{(r)}} x_{ijk} &\leq 1
    \quad\forall i\in\mathcal{R}^{(r)}
    \tag{T2: Single assignment}
\end{align}
$\mathcal{R}^{(r)}$ = tập yêu cầu đặt xe hạng $r$.

\textbf{Giải bằng $\varepsilon$-constraint / Bi-level:}
\begin{enumerate}
  \item Giải bài toán \emph{restricted master} cho hạng Premium trước
    (fixed capacity $|\mathcal{K}^{(3)}|$).
  \item Thêm ràng buộc $W^{(3)}\leq W^{(3)*}+\delta$ vào master.
  \item Giải tiếp cho hạng Plus, rồi Standard.
\end{enumerate}
\end{ridebox}

\textbf{SLA Guarantee:} Khách Premium được cam kết ETA tối đa $T_{\max}^{(3)}=5$ phút;
ràng buộc cứng này được đưa vào bài toán assignment như:
\begin{equation}
  x_{ijk}=0 \quad\text{nếu}\; \hat{\tau}(\ell_k,pu_j) > T_{\max}^{(r_j)}
\end{equation}
tức là không ghép cặp nếu xe không thể đến kịp trong SLA.

%% ================================================================
\section{Multi-Agent RL cho Hệ thống Toàn cục}

\subsection{CTDE --- Centralized Training, Decentralized Execution}

\begin{formulabox}{QMIX --- Monotonic Value Function Factorisation}
\begin{align}
  \pi_k^*(a_k|o_k) &= \argmax_{a_k} Q^{\mathrm{mix}}_\theta(\bm{o},\bm{a})\\
  Q^{\mathrm{mix}}_\theta(\bm{o},\bm{a}) &= f_\theta(Q_1(o_1,a_1),\ldots,Q_m(o_m,a_m))
\end{align}
\end{formulabox}

Đảm bảo tính chất \textbf{IGM} (Individual-Global Max): quyết định tham lam
của từng agent suy ra từ global optimum.

Trạng thái $o_k$ của agent $k$ (xe điện Green SM) bổ sung:
$o_k = (\ell_k, e_k, \text{orders assigned}, \text{nearest charger ETA})$.

\subsection{Continuous Learning và Feedback Loop}

\begin{equation}
  \theta_{t+1}
  = \theta_t - \eta\,\nabla_\theta\,
    \mathcal{L}(\theta_t;\,\mathcal{D}_{\mathrm{replay}})
\end{equation}

Dữ liệu thực tế (SoC thực tế vs dự báo, hủy chuyến, thời gian chờ sạc)
liên tục cập nhật vào \emph{Prioritized Experience Replay}.
Online fine-tuning mỗi giờ bằng \emph{Curriculum Learning}
(đơn giản $\rightarrow$ phức tạp).

%% ================================================================
\section{Tóm tắt Pipeline \& Hướng Phát triển}

\subsection{Bản đồ từ bài toán đến kỹ thuật (cập nhật)}

\begin{table}[h!]
\centering
\caption{Mapping đầy đủ: Bài toán -- Mô hình -- Kỹ thuật DS/AI}
\renewcommand{\arraystretch}{1.25}
\begin{tabular}{p{3.4cm}p{3.6cm}p{4.2cm}}
\toprule
\textbf{Bài toán} & \textbf{Công thức} & \textbf{Kỹ thuật DS/AI} \\
\midrule
\multicolumn{3}{l}{\textit{--- Logistics / Giao hàng ---}}\\
Ghép đơn cùng tuyến & ILP/McCormick & ST-DBSCAN, Greedy Insertion \\
Định tuyến động & DVRP/VRPTW+MTZ & PPO+GAE, Attention Model \\
Dự đoán ETA & Graph Regression & ST-GNN, Quantile+Huber \\
Dự báo cầu & Spatial Time Series & ConvLSTM + GAT \\
Ghép cặp tài xế & Assignment (imbal.) & Hungarian/Auction/$O(n^3)$ \\
Tái phân bổ & Time-Expanded LP & LP / Earth Mover's Distance \\
Robust scheduling & DRO / CVaR & Wasserstein Ball, CVaR-Opt \\
\midrule
\multicolumn{3}{l}{\textit{--- EV-Specific (Green SM) ---}}\\
Routing + sạc điện & E-VRP-CS (SoC) & B\&B, Column Generation \\
Lịch sạc & M/M/c + Reservation & Erlang C, ILP slot booking \\
Hao mòn pin & Lifecycle cost & Rainflow + Arrhenius model \\
Range anxiety & Risk-averse MDP & CVaR-PPO, Reward Shaping \\
\midrule
\multicolumn{3}{l}{\textit{--- Ride-hailing (Chở khách) ---}}\\
Ghép chuyến pool & Shareability Network & Max Weighted Matching \\
Định tuyến khách & DARP (Cordeau) & B\&P, LNS heuristic \\
Hủy / từ chối & Stochastic Matching & Logistic Reg, Online BM \\
Surge pricing & Cân bằng cung-cầu & Causal Inference, A/B test \\
Điểm đón chính xác & POI Clustering & DBSCAN + Map Embedding \\
\midrule
\multicolumn{3}{l}{\textit{--- Hệ thống toàn cục ---}}\\
Phối hợp đa agent & MARL (MDP) & QMIX, CTDE \\
\bottomrule
\end{tabular}
\end{table}

\subsection{Target Latency theo Module (cập nhật)}

\begin{table}[h!]
\centering
\caption{Latency và complexity của từng module}
\renewcommand{\arraystretch}{1.25}
\begin{tabular}{llcc}
\toprule
\textbf{Module} & \textbf{Mô hình} & \textbf{Complexity} & \textbf{Latency} \\
\midrule
Demand Forecast & ConvLSTM+GAT & $O(T|\mathcal{V}|^2)$ & $\sim$100ms \\
ETA Prediction & ST-GNN+Quantile & $O(L|\mathcal{E}|)$ & $\sim$20ms \\
Order Batching & ST-DBSCAN & $O(n^2)$ & $\sim$50ms \\
Ride Pooling & Shareability Graph & $O(n^2)$+Matching & $\sim$80ms \\
Repositioning & Time-Expanded LP & $O(n^3)$ & $\sim$500ms \\
Dispatch & Hungarian/Auction & $O(n^3)/O(n^2\log n)$ & $\sim$10ms \\
DARP Route & B\&P heuristic & NP-hard & $\sim$300ms \\
E-VRP-CS & B\&B+SoC & NP-hard & $\sim$500ms \\
Charging Schedule & M/M/c+ILP & $O(|\mathcal{F}|\cdot T)$ & $\sim$200ms \\
Surge Pricing & Equilibrium model & $O(|G|)$ & $\sim$5ms \\
MARL (QMIX) & Decentralized PPO & $O(m\cdot|\mathcal{A}|)$ & $\sim$1ms \\
\bottomrule
\end{tabular}
\end{table}

\subsection{Hướng Nghiên cứu Tiếp theo}

\begin{enumerate}[leftmargin=2em, itemsep=4pt]
  \item \textbf{V2G (Vehicle-to-Grid):} xe điện rỗi bán điện lại lưới,
        tối ưu chi phí năng lượng và tạo nguồn thu phụ.
  \item \textbf{Multi-modal Trips:} phối hợp Xanh SM Bike $\to$ Metro $\to$ Taxi,
        tối ưu chuyến đi đa phương tiện.
  \item \textbf{LLM + Optimization:} dùng LLM làm planner, giao solver chuyên biệt.
  \item \textbf{Foundation Models cho Time Series:} TimesFM, Chronos
        cho dự báo cầu zero-shot.
  \item \textbf{Diffusion Models cho Route Generation:}
        sinh lộ trình bằng score-based generative models.
  \item \textbf{Federated Learning:} học mô hình phân tán từ nhiều thành phố
        mà không chia sẻ dữ liệu nhạy cảm.
\end{enumerate}

\subsection{Tài liệu Tham khảo}

\begin{enumerate}[label={[\arabic*]}, leftmargin=3em, itemsep=2pt]
  % Routing & VRP
  \item Kool, W.\ et al.\ (2019). \textit{Attention, Learn to Solve Routing Problems.}
        ICLR 2019.
  \item Nazari, M.\ et al.\ (2018). \textit{RL for Solving the VRP.} NeurIPS 2018.
  \item Li, Y.\ et al.\ (2021). \textit{Learning to Delegate for Large-Scale VRP.}
        NeurIPS 2021.
  % RL
  \item Schulman, J.\ et al.\ (2017). \textit{PPO Algorithms.} arXiv:1707.06347.
  \item Rashid, T.\ et al.\ (2018). \textit{QMIX: Monotonic Value Function
        Factorisation for MARL.} ICML 2018.
  % GNN
  \item Veličković, P.\ et al.\ (2018). \textit{Graph Attention Networks.} ICLR 2018.
  \item Shi, X.\ et al.\ (2015). \textit{Convolutional LSTM Network.} NeurIPS 2015.
  % EV
  \item Schneider, M., Stenger, A.\ \& Goeke, D.\ (2014).
        \textit{The E-VRPTW with Recharging Stations.} Transportation Science.
  \item Sweda, T., Dolinskaya, I.\ \& Klabjan, D.\ (2017).
        \textit{Adaptive Routing and Recharging Policies for EVs.}
        Transportation Science.
  \item Keskin, M.\ \& Çatay, B.\ (2016).
        \textit{Partial Recharge Strategies for the E-VRPTW.}
        Computers \& Operations Research.
  \item Pelletier, S., Jabali, O.\ \& Laporte, G.\ (2016).
        \textit{50 Years of Vehicle Routing and EV Routing: Review.}
        Transportation Research Part B.
  % Ride-hailing
  \item Cordeau, J.-F.\ \& Laporte, G.\ (2007).
        \textit{The Dial-a-Ride Problem.} 4OR.
  \item Santi, P.\ et al.\ (2014).
        \textit{Quantifying the Benefits of Vehicle Pooling with
        Shareability Networks.} PNAS.
  \item Alonso-Mora, J.\ et al.\ (2017).
        \textit{On-Demand High-Capacity Ride-Sharing via Dynamic
        Trip-Vehicle Assignment.} PNAS.
  \item Castillo, J., Knoepfle, D.\ \& Weyl, G.\ (2017).
        \textit{Surge Pricing Solves the Wild Goose Chase.}
        ACM EC 2017.
  \item Özkan, E.\ \& Ward, A.\ (2020).
        \textit{Dynamic Matching for Real-Time Ride-Sharing.}
        Stochastic Systems.
  % Robust / Stochastic
  \item Mohajerin Esfahani, P.\ \& Kuhn, D.\ (2018).
        \textit{Data-Driven DRO using Wasserstein Metric.} Math. Programming.
  \item Rockafellar, R.T.\ \& Uryasev, S.\ (2000).
        \textit{Optimization of Conditional Value-at-Risk.}
        Journal of Risk.
  % Hungarian / Matching
  \item Bertsekas, D.P.\ (1988).
        \textit{The Auction Algorithm: A Distributed Relaxation Method
        for the Assignment Problem.} Annals of Operations Research.
\end{enumerate}

\end{document}
